\documentclass[11pt,a4paper]{article}

\usepackage[utf8]{inputenc}
\usepackage[T1]{fontenc}
\usepackage[italian]{babel}
\usepackage{lmodern}
\usepackage{microtype}
\usepackage[a4paper,margin=24mm]{geometry}
\usepackage{xcolor}
\usepackage{booktabs}
\usepackage{tabularx}
\usepackage{longtable}
\usepackage{enumitem}
\usepackage{hyperref}
\usepackage{fancyhdr}
\usepackage{titlesec}

\definecolor{ReferenceOrange}{HTML}{FF9F43}
\definecolor{MixBlue}{HTML}{59BFFF}
\definecolor{DarkPanel}{HTML}{151D27}
\definecolor{LinkBlue}{HTML}{2474A6}

\hypersetup{
  colorlinks=true,
  linkcolor=LinkBlue,
  urlcolor=LinkBlue,
  pdftitle={ReferenceLab 2.1 -- Manuale d'uso},
  pdfauthor={ReferenceLab}
}

\setlist{nosep,leftmargin=6mm}
\setlength{\parindent}{0pt}
\setlength{\parskip}{5pt}
\setcounter{tocdepth}{2}
\pagestyle{fancy}
\fancyhf{}
\lhead{ReferenceLab 2.1}
\rhead{Manuale d'uso}
\cfoot{\thepage}

\newcommand{\control}[1]{\textbf{\textsf{#1}}}
\newcommand{\menuitem}[1]{\textsf{#1}}
\newcommand{\mix}{\textcolor{MixBlue}{\textbf{Mix}}}
\newcommand{\reference}{\textcolor{ReferenceOrange}{\textbf{Reference}}}
\newcommand{\important}[1]{\fcolorbox{ReferenceOrange}{white}{\parbox{0.93\linewidth}{#1}}}

\title{\vspace{-15mm}\textbf{ReferenceLab 2.1}\\[2mm]
\Large Manuale completo di utilizzo}
\author{Plugin VST3 e applicazione Standalone}
\date{Luglio 2026}

\begin{document}
\maketitle

\begin{center}
\small Documento ricavato da ReferenceLab SRS 2.1, cronologia del progetto e
comportamento implementato nel codice sorgente.
\end{center}

\tableofcontents
\newpage

\section{Introduzione}

ReferenceLab è uno strumento di monitoraggio e confronto A/B per confrontare il
segnale della DAW, chiamato \mix, con una traccia commerciale o tecnica,
chiamata \reference. Il plugin riunisce riproduzione delle reference, loudness
matching, filtri di confronto, analisi spettrale e dinamica, librerie, sorgenti
web e correzione cuffie AutoEQ.

La versione 2.1 è disponibile nei formati \textbf{VST3} e
\textbf{Standalone}, opera in stereo e si rivolge principalmente a Windows
10/11 x64. Nel formato VST3 va inserita sul bus sul quale si desidera effettuare
il controllo, normalmente il master bus. Le funzioni del plugin sono pensate
per il monitoraggio: la modalità \control{Safe Export}, attiva inizialmente,
protegge il rendering offline dalle elaborazioni di confronto.

\subsection{Concetti fondamentali}

\begin{description}[style=nextline]
  \item[\mix] Il segnale ricevuto in ingresso dalla DAW.
  \item[\reference] Il file locale o remoto scelto come termine di paragone.
  \item[A/B] Il passaggio tra Mix e Reference. Il comando è smussato dal tempo
  \control{A/B Fade} per evitare click.
  \item[Loudness matching] Compensazione del volume percepito. Serve a evitare
  che il segnale più forte sembri automaticamente migliore.
  \item[Processing di confronto] Filtri HP, banda parametrica, LP e modalità di
  ascolto applicati in modo coerente al confronto.
  \item[Correzione cuffie] EQ finale dedicato al sistema di ascolto. È distinto
  dall'EQ di confronto e viene applicato soltanto se non è in bypass.
\end{description}

\important{\textbf{Attenzione al livello.} Prima di confrontare due sorgenti,
attivare il loudness matching oppure regolare il guadagno manuale. Cambiamenti
di livello possono essere scambiati facilmente per cambiamenti timbrici.}

\section{Avvio rapido}

\begin{enumerate}
  \item Inserire ReferenceLab sul master della DAW oppure avviare la versione
  Standalone.
  \item Aprire \control{FILE} e scegliere \menuitem{Add Audio File...}; in
  alternativa trascinare/aggiungere una sorgente web con
  \menuitem{Add Web Source...}.
  \item Selezionare una reference dalla libreria. La traccia viene caricata e
  avviata.
  \item Usare i pulsanti \control{MIX} e \control{REFERENCE}, oppure il tasto
  \control{B}, per alternare le sorgenti.
  \item Nella pagina \control{Compare}, togliere il bypass del loudness match,
  scegliere il detector e attivare \control{Auto Match}.
  \item Usare \control{Analysis} per verificare spettro, scope, LUFS, picchi,
  dinamica e immagine stereo.
  \item Lasciare \control{Safe Export} attivo se ReferenceLab rimane sul master
  durante l'esportazione offline.
\end{enumerate}

\section{Interfaccia generale}

\subsection{Barra superiore}

La barra superiore contiene:

\begin{longtable}{p{0.24\linewidth}p{0.70\linewidth}}
\toprule
\textbf{Controllo} & \textbf{Funzione}\\
\midrule
\endhead
\control{LIBRARY} & Apre o chiude la libreria laterale. Il contrasto
chiaro/scuro indica lo stato.\\
\control{FILE} & Aggiunta, scansione, sorgenti web, cataloghi JSON,
ricollegamento e rimozione delle reference.\\
\control{CATALOG} & Gestione di librerie, playlist, metadati e preset di
confronto.\\
\control{MIX}/\control{REFERENCE} & Seleziona la sorgente ascoltata. Il simbolo
``\textgreater'' e il colore mostrano la sorgente attiva.\\
\control{B} & Scorciatoia da tastiera per alternare Mix e Reference.\\
\control{\(\leftarrow\)}/\control{\(\rightarrow\)} & Carica la reference
precedente o successiva nell'elenco filtrato.\\
\control{Ingranaggio} & Apre o chiude le impostazioni. Quando è aperto il
pulsante diventa chiaro con simbolo scuro.\\
\bottomrule
\end{longtable}

La riga di stato indica la sorgente ascoltata, riproduzione, pausa, caricamento
web, decodifica, errori e sincronizzazione DAW. La riga host mostra, quando
forniti dalla DAW, play/stop, BPM, metrica e posizione in secondi.

\subsection{Guida contestuale}

La zona centrale che normalmente mostra \control{ReferenceLab 2.1} funziona
anche come guida rapida. Lasciando fermo il puntatore sopra un controllo per
circa 700 ms, il titolo viene sostituito da una breve frase inglese che ne
descrive la funzione. Per esempio, sopra \control{Auto Match} appare
\texttt{Match Mix and Reference loudness.}

Spostando il puntatore su un altro controllo, il tempo di attesa riparte.
Uscendo dal controllo o passando sopra una zona senza descrizione ricompare
\control{ReferenceLab 2.1}. La guida non richiede clic, non modifica valori e
non impedisce trascinamento, rotella o uso normale dell'interfaccia. I testi
sono intenzionalmente brevi per rimanere leggibili nella barra superiore.
Non vengono aperti tooltip o messaggi flottanti: la descrizione appare
esclusivamente al posto del titolo e non copre controlli, grafici o cursori.

\subsection{Pagine}

Le cinque pagine principali sono:

\begin{itemize}
  \item \control{Reference}: player, waveform, loop e sincronizzazione;
  \item \control{Analysis}: spettro, oscilloscopio e misure tecniche;
  \item \control{Compare}: loudness matching, EQ di confronto e campo stereo;
  \item \control{Meter}: storico Airwindows di loudness, transienti e densità;
  \item \control{Headphones}: ricerca AutoEQ e correzione cuffie.
\end{itemize}

\section{Libreria delle reference}

\subsection{Apertura, dimensione e selezione}

La libreria è un pannello laterale richiamato con \control{LIBRARY}. Il
separatore verticale permette di modificarne la larghezza; un doppio clic sul
separatore ripristina la proporzione predefinita. La larghezza e lo stato
aperto/chiuso vengono memorizzati.

Selezionando una riga si carica la reference, si aggiornano waveform e metadati
e si rende la traccia disponibile per il confronto. Le frecce della barra
superiore navigano nell'elenco attualmente visibile.

\subsection{Ricerca, filtri e ordinamento}

La ricerca considera titolo, artista, album, genere e tag. Sono disponibili
anche filtri dedicati per artista, genere, tonalità, BPM minimo/massimo,
preferiti e valutazione minima.

L'ordinamento può usare titolo, artista, album, genere, anno, BPM, rating, data
di aggiunta o ultimo utilizzo. Il comando discendente inverte l'ordine.
Ricerca, filtro e ordinamento cambiano la vista: non eliminano file né
metadati.

\subsection{Menu FILE}

\begin{description}[style=nextline]
  \item[\menuitem{Add Audio File...}] Aggiunge un singolo file audio locale.
  Il progetto usa i decoder registrati da JUCE; i formati effettivamente
  disponibili dipendono dalla build e dalla piattaforma.
  \item[\menuitem{Scan Folder...}] Analizza una cartella e aggiunge i file
  audio rilevati. L'opzione \control{Recursive Scan} include le sottocartelle.
  \item[\menuitem{Add Web Source...}] Accetta un URL HTTP/HTTPS audio. Il
  contenuto viene scaricato e decodificato in memoria; lo stato mostra
  percentuale, decodifica o errore.
  \item[\menuitem{Import JSON Catalog...}] Importa un catalogo JSON locale.
  \item[\menuitem{Relink Reference...}] Associa di nuovo una voce a un file
  locale spostato o rinominato, conservandone i metadati.
  \item[\menuitem{Remove Reference}] Rimuove la voce selezionata dalla libreria.
  Non va usato come comando di cancellazione del file audio originale.
\end{description}

\subsection{Librerie e playlist}

\control{New Library...} crea una libreria associata a una cartella scelta
dall'utente. Il selettore libreria passa da un database all'altro.
\control{New Playlist...} crea una raccolta logica; \control{Add Selection to
Playlist} vi aggiunge la reference selezionata. Le playlist non duplicano i
file audio.

\subsection{Metadati}

Nel pannello impostazioni si possono modificare:

\begin{itemize}
  \item titolo, artista, album e genere;
  \item anno, BPM e tonalità;
  \item rating da 0 a 5 e stato preferito;
  \item tag separati da virgola e note libere;
  \item offset iniziale e punti del loop, tramite la pagina Reference.
\end{itemize}

\control{Save Metadata} salva le modifiche nel database della libreria. È
consigliabile salvare prima di cambiare reference quando sono stati modificati
campi testuali o marker.

\section{Pagina Reference}

\subsection{Waveform}

La waveform rappresenta l'intera reference e visualizza posizione corrente,
inizio e fine del loop. La regione del loop è evidenziata.

\begin{itemize}
  \item Clic o trascinamento in un'area libera: seek.
  \item Trascinamento dei marker: modifica Loop A o Loop B.
  \item Trascinamento per panoramica nella vista ingrandita: sposta la finestra
  temporale.
  \item Rotella del mouse: zoom orizzontale, centrato sulla posizione del
  puntatore.
  \item \control{Full View}: ripristina la visualizzazione dell'intero file.
\end{itemize}

\subsection{Transport}

\control{Play} avvia la reference; \control{Pause} conserva la posizione;
\control{Stop} interrompe e riporta il player alla posizione iniziale prevista.
\control{Position} permette il seek numerico/lineare.

\subsection{Offset e loop}

\control{Start Offset} stabilisce il punto iniziale della reference.
\control{Loop} abilita la ripetizione tra \control{Loop A} e \control{Loop B}.
\control{Full Loop} estende l'intervallo alla durata disponibile. I valori
possono essere regolati con i controlli oppure trascinando i marker.

\subsection{Sincronizzazione DAW}

\control{Transport Sync} collega il comportamento del player al transport host.
Sono disponibili:

\begin{description}[style=nextline]
  \item[\control{Timeline}] La posizione della reference segue la timeline
  comunicata dalla DAW, tenendo conto dell'offset.
  \item[\control{Restart on Play}] La reference riparte quando il transport
  passa allo stato play.
  \item[\control{Sync Offset}] Anticipa o ritarda la relazione con la DAW, in
  secondi, nell'intervallo da -600 a +600.
\end{description}

Se l'host non comunica una posizione valida, l'interfaccia mostra
\control{SYNC N/A}; in questo caso usare la modalità libera.

\section{Confronto A/B e routing}

Il passaggio \control{MIX}/\control{REFERENCE} avviene con crossfade.
\control{A/B Fade} imposta la durata fra 0 e 100 ms. Valori brevi rendono il
confronto immediato; valori più lunghi riducono maggiormente click e discontinuità.

La Reference deve essere caricata perché il relativo pulsante produca audio.
Lo stato \control{REF | NOT LOADED} segnala che non è disponibile.

\section{Pagina Analysis}

\subsection{Spettro in tempo reale}

Lo spettro mostra Mix e Reference con colori coerenti in tutta l'interfaccia.
La sorgente ascoltata ha anche un'area riempita, l'altra conserva la sagoma.
L'asse delle frequenze è logaritmico da circa 20 Hz a 20 kHz; l'asse verticale
è espresso in dB.

\control{Listen Mode} determina sia ciò che si ascolta sia il contenuto mostrato
nello spettro principale:

\begin{description}[style=nextline]
  \item[\control{Stereo}] Normale ascolto stereo; una curva per sorgente.
  \item[\control{Mono}] Somma coerente dei canali per controllare compatibilità
  mono e cancellazioni.
  \item[\control{Mid}] Componente comune \(M=(L+R)/2\).
  \item[\control{Side}] Componente differenziale \(S=(L-R)/2\), utile per
  riverberi, ampiezza e contenuto esclusivamente laterale.
\end{description}

\control{Swap L/R} scambia canale sinistro e destro per entrambe le sorgenti.
Il risultato è rappresentato dallo scope e dai meter Output; non cambia gli
spettri comparativi Mid/Side della pagina Compare.

\subsection{FFT Size, Refresh, Smoothing e Slope}

\begin{description}[style=nextline]
  \item[\control{FFT Size}] Da 1024 a 32768 campioni. Una FFT grande migliora
  la risoluzione alle basse frequenze ma reagisce più lentamente; una piccola
  è più rapida sui transienti.
  \item[\control{Refresh}] Aggiornamento grafico a 15, 30 o 60 Hz. Non modifica
  l'audio.
  \item[\control{Smoothing}] Stabilizza visivamente lo spettro. Valori maggiori
  producono una lettura più lenta e uniforme.
  \item[\control{Slope}] Compensazione esclusivamente visiva da 0,0 a
  4,5 dB/ottava, predefinita a 4,5. Non equalizza il segnale e non modifica i
  dati FFT originali.
\end{description}

\control{Freeze} congela la rappresentazione per effettuare una lettura;
modificare Slope continua comunque ad aggiornare il modo in cui i dati
congelati vengono disegnati. \control{Reset Meters} azzera le misure accumulate.

\subsection{Oscilloscopio Output}

Lo scope mostra il segnale effettivamente instradato dopo modalità di ascolto.
La rotella effettua zoom orizzontale attorno al puntatore. È possibile
allargare la finestra anche sotto il 100\%, per osservare un intervallo più
lungo, oppure superare il 100\% per analizzare il dettaglio. Il trascinamento
nella vista ingrandita sposta la finestra temporale.

\subsection{Meter tecnici}

\begin{longtable}{p{0.27\linewidth}p{0.67\linewidth}}
\toprule
\textbf{Misura} & \textbf{Interpretazione}\\
\midrule
\endhead
Momentary LUFS & Loudness su finestra molto breve; reagisce rapidamente.\\
Short-Term LUFS & Loudness su alcuni secondi; utile per sezioni musicali.\\
Integrated LUFS & Media gated accumulata; va azzerata quando inizia una nuova
misurazione.\\
Peak & Picco campione rilevato.\\
True Peak & Stima dei picchi tra campioni, utile per valutare margine e
conversioni.\\
RMS & Energia media del segnale.\\
Crest Factor & Differenza tra picco e livello medio; indica il rapporto fra
transienti e corpo.\\
Dynamic Range & Stima dell'escursione dinamica.\\
Correlation & Relazione di fase L/R: valori positivi sono più compatibili mono,
valori negativi richiedono attenzione.\\
Stereo Width & Stima normalizzata della larghezza stereo.\\
\bottomrule
\end{longtable}

\section{Pagina Compare}

\subsection{Loudness Match}

\control{Match Bypass} esclude l'intera compensazione. Con bypass disattivato:

\begin{itemize}
  \item \control{Auto Match} calcola la differenza fra Mix e Reference;
  \item \control{Detector} sceglie Integrated, Short-Term o Momentary;
  \item \control{Target LUFS} stabilisce il riferimento operativo, da -36 a
  -6 LUFS;
  \item \control{Manual Gain} aggiunge una correzione manuale da -12 a
  +12 dB;
  \item \control{Applied} mostra il guadagno corrente applicato.
\end{itemize}

La compensazione automatica è limitata per evitare guadagni estremi ed è
smussata nel tempo. Integrated è più stabile per confronti lunghi; Short-Term
è adatto alle sezioni; Momentary è rapido ma può muoversi di più.

\subsection{Spettri post-EQ}

Il grafico mostra gli spettri post-EQ Mid e Side di Mix e Reference. I pulsanti
\control{MID} e \control{SIDE} abilitano i livelli indipendentemente dalla
modalità Listen della pagina Analysis. Disattivandoli entrambi non viene
mostrato alcuno spettro.

La sorgente selezionata è riempita, quella non selezionata resta come contorno;
il Side è reso più trasparente del Mid. La curva bianca mostra la risposta
complessiva dei filtri di confronto quando l'EQ non è in bypass.

\subsection{Campo stereo comparativo}

Il visualizzatore inferiore rappresenta contemporaneamente il campo stereo di
Mix e Reference. L'asse orizzontale evidenzia la componente laterale, quello
verticale la componente Mid. Il riempimento identifica la sorgente ascoltata.
Il grafico è diagnostico e non modifica l'audio.

\subsection{Comparison Filters}

\control{EQ Bypass} esclude in blocco i filtri seguenti:

\begin{description}[style=nextline]
  \item[\control{HP Enabled / High Pass}] Passa-alto da 20 a 2000 Hz con
  pendenza 12 o 24 dB/ottava.
  \item[\control{Band Enabled / Band Freq / Gain / Q}] Banda parametrica da
  20 Hz a 20 kHz, guadagno da -18 a +18 dB e Q da 0,1 a 12. Un Q alto
  seleziona una zona stretta; un Q basso una zona ampia.
  \item[\control{LP Enabled / Low Pass}] Passa-basso da 100 Hz a 20 kHz con
  pendenza 12 o 24 dB/ottava.
\end{description}

I filtri vengono applicati in modo equivalente a Mix e Reference: servono a
confrontare una regione dello spettro, non a ``migliorare'' una sola sorgente.

\section{Pagina Meter}

Il meter ispirato ad Airwindows Meter registra una timeline separata per Mix e
Reference. Le tre corsie sono:

\begin{description}[style=nextline]
  \item[Loudness / Peak] Mostra livello medio e picchi per L/R.
  \item[Detail / Slew] Evidenzia velocità dei cambiamenti e contenuto
  transiente.
  \item[Fullness / Zero Cross] Usa lo zero crossing come indicazione della
  distribuzione verso basse o alte frequenze.
\end{description}

I riquadri MIX e REF producono valutazioni cumulative per le tre corsie. Sono
indicatori comparativi, non standard normativi di loudness. Il colore dei punti
comunica l'attività relativa di picco e slew; i colori viola/turchese
distinguono i canali nella corsia zero crossing. Le linee verticali bianche
segnano i cambi sorgente.

Un doppio clic sul meter cancella storico e punteggi. Per un confronto utile,
ascoltare porzioni di durata simile e azzerare prima di una nuova sessione.

\section{Pagina Headphones e AutoEQ}

\subsection{Scopo e posizione nella catena}

La correzione cuffie è un EQ finale di monitoraggio. È separata dai filtri
Compare: questi ultimi isolano frequenze su entrambe le sorgenti, mentre
Headphones compensa il modello di cuffia usato per ascoltare l'uscita finale.

\important{\textbf{Un solo profilo alla volta.} I profili salvati non vengono
sommati. Il doppio clic carica un singolo profilo attivo, indicato dal simbolo
``\textgreater''. \control{BYPASS ON} significa che nessuna correzione cuffie
viene applicata.}

\subsection{Ricerca online}

La sezione \control{AUTOEQ SEARCH} è compatta per impostazione predefinita.
Aprirla, digitare almeno due caratteri del modello e premere
\control{SEARCH AUTOEQ}. La ricerca richiede Internet.

I risultati indicano modello, sorgente della misura e rig quando disponibili.
Selezionare la misura desiderata. Il campo nome è opzionale: se rimane vuoto,
il profilo viene salvato con il nome originale della cuffia. Premere
\control{SAVE ONLINE PROFILE} per richiedere/generare l'EQ, salvarlo localmente
e caricarlo.

Dopo il salvataggio il profilo è utilizzabile offline: non è necessario
interrogare nuovamente il servizio finché non si vuole cercare una nuova
cuffia.

\subsection{Profili locali}

\control{PROFILES} apre o chiude la barra laterale. Un clic seleziona il profilo
per visualizzarlo e modificarlo; un \textbf{doppio clic} lo carica nel DSP. Il
simbolo ``\textgreater'' identifica quello realmente attivo.

\control{SAVE CHANGES} salva nome, safety preamp e bande modificate. Se il
profilo modificato è attivo, il DSP viene aggiornato. \control{DELETE PROFILE}
elimina il profilo locale; se era attivo, la correzione viene rimossa.

\subsection{Bypass}

Il bypass è attivo inizialmente:

\begin{itemize}
  \item \control{BYPASS ON}: correzione disabilitata;
  \item \control{BYPASS OFF}: profilo attivo applicato all'audio.
\end{itemize}

Prima di togliere il bypass abbassare il volume: alcuni profili modificano
significativamente l'equilibrio tonale, anche se il preamp negativo fornisce
margine.

\subsection{Safety Preamp e bande}

\control{EQ BANDS} espande o compatta l'editor. \control{Safety Preamp}, da
-40 a 0 dB, precede i filtri e riduce il rischio di clipping dovuto ai boost.

Ogni banda contiene:

\begin{itemize}
  \item tipo \control{Peaking}, \control{Low Shelf} o \control{High Shelf};
  \item frequenza da 10 a 40000 Hz;
  \item gain da -30 a +30 dB;
  \item Q da 0,05 a 30.
\end{itemize}

Il grafico bianco mostra la risposta totale, preamp incluso, e si aggiorna
mentre si muovono i controlli. Le modifiche diventano permanenti solo dopo
\control{SAVE CHANGES}.

Safety Preamp e bande si trovano nella stessa area scorrevole. Per evitare che
la regolazione dei knob sposti il pannello, lo scorrimento si effettua
trascinando esclusivamente la barra laterale.

\subsection{Profili simili e cambio profilo}

Due risultati con lo stesso modello possono provenire da sorgenti, rig o target
diversi e quindi non essere equivalenti. Confrontare sempre preamp e valori
delle bande, non soltanto il nome. Il caricamento ricrea lo stato dei filtri:
anche fra configurazioni identiche può essere percepito un breve transiente nel
momento del cambio; non significa che i profili vengano concatenati.

\section{Impostazioni}

\subsection{Cache}

Il limite della cache RAM può essere 128, 256, 512 o 1024 MB. La cache contiene
audio locale decodificato per rendere più rapidi i cambi. \control{Clear Cache}
libera la RAM senza interrompere la reference già in riproduzione. La cache non
è la libreria e cancellarla non rimuove file o metadati.

\subsection{Aspetto}

I colori Mix e Reference possono essere scelti fra blu, arancione, verde,
viola, giallo e rosso. Il bianco è riservato ai controlli neutrali, alla curva
EQ e ai visualizzatori generali. Colore, etichette, riempimento e contorni
lavorano insieme affinché la sorgente non sia identificata soltanto dal colore.

\subsection{Preset e stato}

\control{Save Comparison Preset...} crea un file \texttt{.rlpreset} con i
parametri di confronto. \control{Load Comparison Preset...} li ripristina.
Questi preset non incorporano il file audio della reference né sostituiscono
la libreria.

Lo stato del plugin salvato dalla DAW comprende parametri, sorgente attiva,
posizione e profilo cuffie attivo. I profili cuffie e il database delle
reference sono inoltre conservati localmente dall'applicazione.

\subsection{Safe Export}

\control{Safe Export} è attivo per impostazione predefinita. Quando l'host
segnala processing non real-time, ReferenceLab non processa il blocco: in
questo modo una reference o una modalità di monitoraggio non finisce
accidentalmente nel bounce. Il comportamento dipende dalla corretta
segnalazione offline della DAW; per sicurezza è comunque buona pratica
disabilitare o rimuovere dal percorso di render i plugin esclusivamente di
monitoraggio.

\section{Workflow consigliati}

\subsection{Confronto timbrico affidabile}

\begin{enumerate}
  \item Caricare una reference dello stesso genere e con caratteristiche
  pertinenti.
  \item Scegliere una sezione confrontabile e impostare il loop.
  \item Attivare Auto Match con detector Short-Term o Integrated.
  \item Alternare con \control{B}, senza guardare inizialmente i grafici.
  \item Usare HP/LP per confrontare separatamente basse, medie e alte.
  \item Consultare gli spettri solo per confermare ciò che si sente.
\end{enumerate}

\subsection{Controllo mono e fase}

\begin{enumerate}
  \item In Analysis scegliere \control{Mono}.
  \item Controllare variazioni di voce, basso e strumenti centrali.
  \item Osservare Correlation: valori persistentemente negativi suggeriscono
  possibili cancellazioni.
  \item Passare a \control{Side} per individuare elementi eccessivamente
  laterali.
  \item Usare il campo stereo Compare per rapportare Mix e Reference.
\end{enumerate}

\subsection{Sessione in cuffia}

\begin{enumerate}
  \item Cercare e salvare il profilo appropriato mentre si è online.
  \item Verificare sorgente, rig e target del risultato.
  \item Fare doppio clic sul profilo e controllare il simbolo ``\textgreater''.
  \item Abbassare il monitor level e passare da \control{BYPASS ON} a
  \control{BYPASS OFF}.
  \item Non compensare nuovamente lo stesso profilo con EQ esterni.
  \item Prima dell'export verificare Safe Export e il routing di monitoraggio.
\end{enumerate}

\section{Risoluzione dei problemi}

\begin{longtable}{p{0.31\linewidth}p{0.63\linewidth}}
\toprule
\textbf{Problema} & \textbf{Controlli consigliati}\\
\midrule
\endhead
La Reference non si sente & Verificare che sia caricata, che il player sia in
Play e che lo stato non mostri NOT LOADED, WEB ERROR o DECODING.\\
SYNC N/A & La DAW non fornisce una posizione utilizzabile; disattivare
Transport Sync o controllare le impostazioni transport dell'host.\\
Il confronto sembra molto diverso & Controllare Match Bypass, Auto Match,
Manual Gain, EQ Bypass, Listen Mode e correzione cuffie. Il livello è spesso la
prima causa.\\
Nessuno spettro in Compare & Riattivare almeno uno fra MID e SIDE.\\
Lo spettro sembra inclinato & Regolare Slope; è una compensazione grafica e non
un EQ.\\
Lo scope è troppo fitto o troppo vuoto & Usare la rotella sopra lo scope; sono
disponibili viste inferiori e superiori al 100\%.\\
La ricerca AutoEQ non risponde & Verificare Internet, usare almeno due
caratteri e riprovare. I profili già salvati restano disponibili offline.\\
Il profilo cuffie non cambia & Il clic singolo mostra soltanto il profilo:
effettuare doppio clic e verificare il simbolo ``\textgreater''. Controllare
anche che BYPASS sia OFF.\\
Due profili omonimi suonano diversi & Controllare source, rig, target, preamp e
tutte le bande: il nome del modello non garantisce parametri identici.\\
I knob cuffie spostano la vista & Usare la barra laterale per lo scroll; il
trascinamento sul knob deve modificare esclusivamente il valore.\\
Il bounce contiene processing indesiderato & Attivare Safe Export e verificare
che la DAW segnali il render offline; in caso contrario bypassare ReferenceLab
prima dell'esportazione.\\
\bottomrule
\end{longtable}

\section{Persistenza, rete e dati}

Le reference locali rimangono nei percorsi originali; ReferenceLab conserva
catalogo e metadati. Le sorgenti web vengono scaricate con limiti di dimensione
e decodificate senza eseguire rete sul thread audio. Cataloghi e URL devono
usare HTTP/HTTPS validi.

I profili cuffie AutoEQ sono salvati nel database locale
\texttt{ReferenceLab/headphones.json} della cartella dati applicazione
dell'utente. Ogni profilo contiene identificativo, nome, modello, sorgente,
forma, rig, target, preamp e fino a 16 filtri. È possibile utilizzarli senza
connessione, modificarli o eliminarli dall'interfaccia.

La disponibilità e i contenuti del servizio AutoEQ dipendono dal servizio
online e dalle rispettive misurazioni. ReferenceLab non sostituisce una
calibrazione individuale della singola unità di cuffie.

\section{Riferimento rapido dei parametri audio}

\begin{longtable}{p{0.28\linewidth}p{0.22\linewidth}p{0.43\linewidth}}
\toprule
\textbf{Parametro} & \textbf{Intervallo/scelte} & \textbf{Effetto}\\
\midrule
\endhead
A/B Fade & 0--100 ms & Crossfade fra Mix e Reference.\\
Target LUFS & -36-- -6 LUFS & Obiettivo del loudness matcher.\\
Manual Gain & -12--+12 dB & Correzione aggiuntiva del match.\\
HPF & 20--2000 Hz & Passa-alto comparativo.\\
HPF Slope & 12/24 dB/ott. & Pendenza passa-alto.\\
Bell Frequency & 20--20000 Hz & Centro della banda parametrica.\\
Bell Gain & -18--+18 dB & Boost/cut della banda.\\
Bell Q & 0,1--12 & Larghezza della banda.\\
LPF & 100--20000 Hz & Passa-basso comparativo.\\
LPF Slope & 12/24 dB/ott. & Pendenza passa-basso.\\
Sync Offset & -600--+600 s & Scostamento dalla timeline DAW.\\
FFT Size & 1024--32768 & Risoluzione/reattività grafica.\\
FFT Slope & 0--4,5 dB/ott. & Compensazione soltanto visiva.\\
Headphone Preamp & -40--0 dB & Margine prima dell'EQ cuffie.\\
Headphone Frequency & 10--40000 Hz & Frequenza della banda cuffie.\\
Headphone Gain & -30--+30 dB & Guadagno della banda cuffie.\\
Headphone Q & 0,05--30 & Larghezza/risonanza della banda cuffie.\\
\bottomrule
\end{longtable}

\section{Note finali}

ReferenceLab è uno strumento di decisione: grafici e voti aiutano a trovare
differenze, ma non definiscono da soli la qualità di un mix. Per ottenere
confronti ripetibili usare porzioni equivalenti, livelli compensati, un sistema
di ascolto noto e una sola correzione cuffie attiva.

\vfill
\begin{center}
\small ReferenceLab 2.1 --- Manuale d'uso --- Revisione iniziale, luglio 2026
\end{center}

\end{document}
